We present a neural network (NN) approach to fit and predict implied volatility surfaces (IVSs).
Atypically to standard NN applications, financial industry practitioners use such models equally to replicate market prices and to value other financial instruments.
In other words, low training losses are as important as generalization capabilities.
Importantly, IVS models need to generate realistic arbitrage-free option prices, meaning that no portfolio can lead to risk-free profits.
We propose an approach guaranteeing the absence of arbitrage opportunities by penalizing the loss using soft constraints.
Furthermore, our method can be combined with standard IVS models in quantitative finance, thus providing a NN-based correction when such models fail at replicating observed market prices.
This lets practitioners use our approach as a plug-in on top of classical methods.
Empirical results show that this approach is particularly useful when only sparse or erroneous data are available.
We also quantify the uncertainty of the model predictions in regions with few or no observations.
We further explore how deeper NNs improve over shallower ones, as well as other properties of the network architecture.
We benchmark our method against standard IVS models.
By evaluating our method on both training sets, and testing sets, namely, we highlight both their capacity to reproduce observed prices and predict new ones.

%As is the case with any domain these days, deep learning aims towards pervading applications in finance.
%Equally, it has been one of the more challenging areas to see an effective adoption, since decisions relate to high impact and high sensitivity (or risk) simultaneously.
%This advocates the need for a principled framework in guiding the integration of neural networks in developing successful and trustworthy applications for finance.
%In this work we focus on the smoothing of the implied volatility surface, however, the same methodology extends to other applications in finance.
%The objective of this paper are two fold: first we introduce the requirements and how to implement a realistic neural net model informed by prior knowledge from financial mathematics, second, we present important practical details on how to properly train and evaluate option pricing models.
%Namely, in an extensive empirical study we highlighting both the capacity to reproduce observed prices and predict new ones, by enforcing well-behaved extrapolation so called asymptotic control.